\documentclass[10pt,a4paper]{amsart}
\usepackage[margin=19mm]{geometry}
\usepackage{amsmath,amssymb,mathtools}
\usepackage[scr=rsfs,cal=euler]{mathalfa}
\usepackage{xfrac}
\usepackage[unicode,hidelinks,bookmarks=false]{hyperref}
\newcommand{\ie}{i.e.\ }
\newcommand{\cf}{cf.\ }
\newcommand{\resp}{respectively\ }
\newcommand{\Subsection}{\subsection}
\providecommand{\nobreakdash}{\nobreak\hbox{-}}
\setlength{\emergencystretch}{2em}
\multlinegap0pt
\usepackage{mathtools}
\usepackage[shortlabels]{enumitem}
\newtheorem{claim}{Claim}
\newcommand{\Orb}{\mathcal{O}}
\newcommand{\Vect}{\mathbf{V}}
\DeclareMathOperator{\colim}{colim}
\renewcommand{\mathcal}{\mathscr}
\title{Unipotent morphisms}
\author{Daniel Bragg, Jack Hall, Siddharth Mathur}
\date{Source: arXiv:2110.15041v2 (CC BY 4.0)}
\begin{document}
\maketitle

\noindent\emph{Source and licence.} Unipotent morphisms. Version arXiv:2110.15041v2 (CC BY 4.0), distributed by arXiv under the Creative Commons Attribution 4.0 International licence, http://creativecommons.org/licenses/by/4.0/, as recorded in the arXiv licence metadata for this exact version and shown on the abstract page https://arxiv.org/abs/2110.15041v2. Full licence notice: \texttt{licenses/arxiv-2110.15041-CC-BY-4.0.txt}.\par\smallskip

\noindent\textit{Editorial scope.} Counted unit: the article's claim TC:schappi:colim (source0.tex lines 518--521). The article's complete original proof of it is delivered with it (source0.tex lines 522--555, one \texttt{proof} environment of 34 lines). No mathematics is written, repaired or re-derived: the statement and the proof are the article's own lines, in the article's own notation, and no retained line is substituted or re-flowed.  No further source-local context is needed: every cross-reference of every retained line resolves inside the two delivered spans.  Nothing was omitted from any retained span, so the package carries no omission record.  No retained line contains a citation, so the wrapper carries no bibliography and no bibliography entry.  No retained line contains a figure environment, an image-inclusion command or drawing code, so no image is delivered and the package declares no asset dependency.  Editorial formatting only, in addition: an AMS \texttt{amsart} preamble that restates the article's own declarations and macros used by the retained lines, namely the environments \texttt{claim} on separate counters, together with the standard packages those lines need.  The retained environments print in this document as Claim 1 in order of appearance, whereas the article prints the same items with the numbering of its own full document.  The complete native source of the article as distributed in the arXiv e-print for this exact version is bundled unchanged as \texttt{sources/117796/source0.tex}.
  \begin{claim}\label{TC:schappi:colim}
    Every pair of objects in $\Vect_G$ has an upper bound and
    $\colim(\mu_G) \simeq G$.
  \end{claim}

    \begin{proof}
      \renewcommand{\qedsymbol}{$\blacksquare$} If $(H_1,\eta_1)$ and
      $(H_2,\eta_2) \in \Vect_G$, then
      $(H_1\oplus H_2,\eta_1\oplus \eta_2) \in \Vect_G$. This proves
      the existence of upper bounds for every pair. We now prove that
      $\colim(\mu_G) \simeq G$. To do this, consider a quasi-coherent
      $\Orb_X$-module $N$ together with $\Orb_X$-module homomorphisms
      $\nu_{(H,\eta)} \colon H \to N$ for every $(H,\eta) \in \Vect_G$
      such that if $h\colon (H,\eta) \to (H',\eta')$ is a morphism in
      $\Vect_G$, then $\nu_{(H,\eta)} = \nu_{(H',\eta')} \circ h$. It
      suffices to show that there is a uniquely induced morphism
      $\nu \colon G \to N$ such that $\nu_{(H,\eta)} = \nu \circ \eta$
      for every $(H,\eta) \in \Vect_G$. Since $X$ has the $\Vect$-resolution
      property, $G$ admits a presentation:
      \[
        \bigoplus_{j\in J} P_j \xrightarrow{\oplus p_j}
          \bigoplus_{(H,\eta)\in \Vect_G} H \xrightarrow{\oplus_{(H,\eta)\in\Vect_G} \eta} G \to 0,
      \]
      where $P_j$ belongs to $\Vect$ for all $j\in J$. Let
      $j\in J$ and note that the morphism
      $p_j \colon P_j \to \oplus_{(H,\eta)\in \Vect_G} H$ factors as
      \[
        P_j \xrightarrow{\tilde{p}_j} \bigoplus_{(H,\eta)\in I_j} H
        \subseteq \bigoplus_{(H,\eta)\in \Vect_G} H,
      \]
      where $I_j \subseteq \Vect_G$ is finite. Let
      $\tilde{Q}_j = \oplus_{(H,\eta)\in I_j} H$ and let
      $\tilde{q}_j = \oplus_{(H,\eta) \in I_j} \eta \colon \tilde{Q}_j \to G$ be the
      resulting morphism. Then $\tilde{q}_j\circ \tilde{p}_j = 0$ and
      so $\tilde{p}_j\colon (P_j,0) \to (\tilde{Q}_j,\tilde{q}_j)$ is a morphism
      in $\Vect_G$. Hence, $\nu_{(\tilde{Q}_j,\tilde{q}_j)} \circ \tilde{p}_j=\nu_{(P_j,0)}=\nu_{(P_j,0)} \circ 0=0$. But this means $(\oplus_{(H,\eta) \in \Vect_G} \nu_{(H,\eta)}) \circ (\oplus_j {p}_j) = 0$. By the universal property of cokernels, there is a unique
      morphism $\nu \colon G \to N$ such that
      $\nu \circ \eta = \nu_{(H,\eta)}$ for all $(H,\eta)\in \Vect_G$. 
  \end{proof}

\end{document}
